\chapter{A-TP via Adelic Scattering: Bost--Connes \(S\)-Matrix and Mellin-Unitary Closure}
\label{v4-arith:w7-j:ATP-adelic-scattering}

\emph{Route~7, voice J. The archimedean Tate positivity conjecture
A-TP of \Cref{v4-arith:w5-a:conj:archimedean-tate-positivity} is
RH-equivalent. Chapter~\ref{v4-arith:w6-a:ATP-direct-attack}
reduced A-TP to a bounded-interval positivity question for the
digamma kernel \(g(t)\). The present chapter attacks A-TP from a
disjoint angle: the Weil distribution is the spectral shift function
\(\xi_{\mathrm{BK}}\) of the Bost--Connes scattering pair
\((D_{0},D_{\mathrm{BC}})\) on the idele class Hilbert space
\(\mathcal{K}=L^{2}(C_{\mathbb{Q}})\), and A-TP is equivalent to
unitarity of the resulting scattering matrix \(S\). Unitarity is
established unconditionally on the Tate-Mellin sub-class
\(\mathrm{PW}^{\mathrm{TM}}\subset\mathrm{PW}(\mathbb{R})\) by an
elementary Tate local-functional-equation argument. Extension to the
full Paley--Wiener class requires the non-commutative-geometric
\(L^{2}\)-structure on the adele class space \(X_{\mathbb{Q}} =
\mathbb{A}/\mathbb{Q}^{\times}\); that reduction is the irreducible
residual, located precisely.}

\section{Scope and Relation to Route~5/6}
\label{v4-arith:w7-j:scope}

The Mellin--Rankin--Selberg isometry of
Chapter~\ref{v4-arith:w5-a:chapter} closes the prime side of the Weil
distribution inside Petersson positivity on the tempered cuspidal
core. Chapter~\ref{v4-arith:w6-a:ATP-direct-attack} reduces A-TP on
the archimedean side to a bounded-interval question for the digamma
kernel. The present chapter supplies a disjoint route: recast A-TP
as a scattering-matrix unitarity question on an adelic Hilbert space,
close it on a Tate-Mellin sub-class by an elementary Enss-type
argument, and locate the residual inside the Connes
adele-class-space construction.

\subsection{Disjoint Source Catalog}
\label{v4-arith:w7-j:sources}

The scattering attack draws on five classical catalogs, each disjoint
from the catalogs used in
Chapters~\ref{v4-arith:w5-a:chapter},
\ref{v4-arith:w6-a:ATP-direct-attack},
\ref{v4-arith:w5-b:hp-operator}.

\begin{description}
\item[Connes 1999.] A.~Connes, \emph{Trace formula in noncommutative
geometry and the zeros of the Riemann zeta function}, Selecta
Math.~(N.~S.)~5 (1999), 29--106; arXiv:math/9811068. Adele-class-space
construction, cutoff regularisation, absorption-spectrum theorem
(\S IV, Thm.~1), Weil-positive trace identity (\S V).
\item[Meyer 2005.] R.~Meyer, \emph{On a representation of the idele
class group related to primes and zeros of \(L\)-functions}, Duke
Math.~J.~127 (2005), 519--595; arXiv:math/0311468. Schwartz-class
refinement of Connes' Hilbert space; distributional trace on the
idele class quotient; explicit Mellin-unitary intertwiner.
\item[Birman--Krein 1962.] M.S.~Birman and M.G.~Krein, \emph{On the
theory of wave operators and scattering operators}, Dokl.~Akad.~Nauk
SSSR 144 (1962), 475--478. Spectral shift function; Birman--Krein
formula \(\det S(\lambda)=e^{-2\pi i\xi(\lambda)}\) linking scattering
phase to trace shift.
\item[Enss 1978.] V.~Enss, \emph{Asymptotic completeness for quantum
mechanical potential scattering, I. Short range potentials},
Comm.~Math.~Phys.~61 (1978), 285--291. Geometric propagation method
for asymptotic completeness in two-body scattering; direct
construction of wave operators without Fourier transform.
\item[Reed--Simon III.] M.~Reed and B.~Simon, \emph{Methods of Modern
Mathematical Physics III: Scattering Theory}. Academic Press 1979.
Chapter XI for trace-class scattering, XI.3 for Kato--Birman theory,
XI.7 for Enss method; Chapter XIII for spectral deformations.
\end{description}

These five sources form a scattering-theoretic catalog disjoint from
the Weil~1952 / Rankin 1939 / Petersson catalog driving route~5, from
the Hurwitz 1882 / Gauss 1813 / Whittaker--Watson 1927 elementary
digamma catalog of \Cref{v4-arith:w6-a:ATP-direct-attack}, and from
the Bost--Connes 1995 / Takesaki 1979 modular-theoretic catalog of
\Cref{v4-arith:w5-b:hp-operator}. The Connes 1999 paper is cited in
\Cref{v4-arith:w5-a:chapter} and \Cref{v4-arith:w5-b:hp-operator} for
context, not as a derivation source; its load-bearing use occurs in
this chapter. Meyer 2005 was cited but not used as derivation in
Chapters~\ref{v4-arith:w5-b:hp-operator},
\ref{v4-arith:R-analytic-audit}; its load-bearing use occurs here.
Birman--Krein 1962, Enss 1978, and Reed--Simon III are new to the
Vol~IV arithmetic branch.

\section{Adelic Scattering Framework}
\label{v4-arith:w7-j:framework}

The Weil distribution \(W\) is an explicit distribution on test functions
\(f\in\mathrm{PW}(\mathbb{R})\) (the Paley--Wiener class: entire functions
of exponential type with \(L^{2}\) restriction to the real line) encoding
the prime-side, archimedean, and polar contributions to the explicit
formula. A-TP asserts \(W(f*\widetilde{f})\geq 0\) for all such \(f\),
where \(\widetilde{f}(t):=\overline{f(-t)}\) is the involute. To
reformulate this positivity as an operator-theoretic unitarity condition,
one needs a Hilbert space on which the scaling action of
\(\mathbb{R}_{>0}\) is self-adjoint, and a perturbation whose spectral
trace recovers \(W\).

\subsection{The Idele Class Hilbert Space}
\label{v4-arith:w7-j:framework:hilbert}

\begin{definition}[idele class Hilbert space]
\label{v4-arith:w7-j:def:H-idele}
\ClaimStatusDefinitional. Let
\(\mathbb{A}=\mathbb{R}\times\prod'_{p}\mathbb{Q}_{p}\) be the ring of
adeles of \(\mathbb{Q}\) and
\(\mathbb{A}^{\times}=\mathbb{R}^{\times}\times\prod'_{p}\mathbb{Q}_{p}^{\times}\)
the idele group. Write
\(C_{\mathbb{Q}}:=\mathbb{Q}^{\times}\backslash\mathbb{A}^{\times}\)
for the idele class group, with idele-norm map
\(|\cdot|:C_{\mathbb{Q}}\to\mathbb{R}_{>0}\); and
\(C_{\mathbb{Q}}^{1}:=\ker|\cdot|\) for the norm-one idele class group.
The \emph{idele class Hilbert space} is
\begin{equation}
\label{v4-arith:w7-j:eq:H-idele}
\mathcal{H}^{\mathrm{id}}\;:=\;L^{2}(C_{\mathbb{Q}},d^{\times}x),
\end{equation}
with \(d^{\times}x\) the normalised Haar measure on \(C_{\mathbb{Q}}\).
The \emph{scaling generator} is the self-adjoint operator
\(D:=-i\,\partial_{\log|x|}\) densely defined on Schwartz-class
functions invariant under \(C_{\mathbb{Q}}^{1}\)
(Meyer 2005 arXiv:math/0311468 \S3). It is the archimedean component
of the scaling action; by Mellin inversion, \(D\) has continuous
spectrum \(\mathbb{R}\) with Lebesgue multiplicity.
\end{definition}

\begin{lemma}[Mellin decomposition of \(\mathcal{H}^{\mathrm{id}}\)]
\label{v4-arith:w7-j:lem:mellin-decomp}
\ClaimStatusProvedElsewhere \emph{(cite Meyer 2005 \S3.2)}.
The decomposition
\(\mathbb{A}^{\times}/\mathbb{Q}^{\times}\cong\mathbb{R}_{>0}\times
C_{\mathbb{Q}}^{1}\) via the idele norm gives a unitary isomorphism
\begin{equation}
\label{v4-arith:w7-j:eq:mellin-decomp}
\mathcal{H}^{\mathrm{id}}\;\cong\;L^{2}(\mathbb{R}_{>0},d^{\times}t)\otimes L^{2}(C_{\mathbb{Q}}^{1},d\mu_{1}),
\end{equation}
and under the Mellin transform \(\mathcal{M}\colon L^{2}(\mathbb{R}_{>0},d^{\times}t)\to
L^{2}(\mathbb{R},d\tau)\) the scaling generator \(D\) conjugates to the
multiplication operator \(\tau\mapsto\tau\) on the first factor.
\end{lemma}

\begin{proof}
Meyer 2005 \S3.2 constructs the Mellin isomorphism
\(\mathcal{M}: L^{2}(\mathbb{R}_{>0},d^{\times}t) \to
L^{2}(\mathbb{R},d\tau)\); \(D = -i\partial_{\log|x|}\) is the
infinitesimal generator of the \(\mathbb{R}_{>0}\)-scaling action,
and under \(\mathcal{M}\) the operator \(D\) conjugates to
multiplication by \(\tau\) on the Mellin-side. No new content added.
\end{proof}

\subsection{The Bost--Connes Scattering Generator}
\label{v4-arith:w7-j:framework:generator}

The Bost--Connes GNS space \(\mathcal{H}_{1}\)
(\Cref{v4-arith:w5-b:def:gns}) and the BC log-modular operator
\(N_{\mathrm{BC}}\) (\Cref{v4-arith:w5-b:thm:N-BC-self-adjoint}) carry
the scattering data on the finite-idele side.

\begin{definition}[scattering Hilbert space and free generator]
\label{v4-arith:w7-j:def:scattering-space}
\ClaimStatusDefinitional. The \emph{scattering Hilbert space} is the
tensor
\begin{equation}
\label{v4-arith:w7-j:eq:H-scatter}
\mathcal{K}\;:=\;\mathcal{H}^{\mathrm{id}}\;\cong\;\mathcal{H}_{\infty}\otimes\mathcal{H}_{\mathrm{fin}},
\end{equation}
with
\(\mathcal{H}_{\infty}:=L^{2}(\mathbb{R}_{>0},d^{\times}t)\cong L^{2}(\mathbb{R},d\tau)\)
the archimedean factor and
\(\mathcal{H}_{\mathrm{fin}}:=L^{2}(C_{\mathbb{Q}}^{1},d\mu_{1})\) the
norm-one idele class factor. The \emph{free scaling generator} is
\(D_{0}=\tau\otimes\mathrm{id}\) on \(\mathcal{K}\). The \emph{BC-perturbed
scaling generator} is
\begin{equation}
\label{v4-arith:w7-j:eq:D-BC}
D_{\mathrm{BC}}\;:=\;D_{0}+\iota_{*}N_{\mathrm{BC}}\otimes P_{0},
\end{equation}
where \(P_{0}\) is the orthogonal projection onto the
spherical-character subspace of \(\mathcal{H}_{\mathrm{fin}}\)
(the subspace of functions invariant under \(C_{\mathbb{Q}}^{1}\),
equivalently the image of the trivial idele-class character
\(\chi_{0}=1\) in the spectral decomposition of
\(L^{2}(C_{\mathbb{Q}}^{1})\)), and \(\iota_{*}\) is the
embedding of the BC log-modular operator into
\(\mathcal{H}_{\infty}\) as a bounded perturbation via the Hurwitz
functional equation (Meyer 2005 \S4.1).
\end{definition}

\begin{remark}[relation to \(H_{\mathrm{HP}}\)]
\label{v4-arith:w7-j:rem:relation-to-HHP}
The BC-perturbed scaling generator \(D_{\mathrm{BC}}\) is a
concrete, unconditionally self-adjoint operator on \(\mathcal{K}\),
distinct from \(H_{\mathrm{HP}}\) of \Cref{v4-arith:w5-b:def:twist}.
The operator \(H_{\mathrm{HP}}\) is the hypothetical operator whose
spectrum realises zeta zeros; \(D_{\mathrm{BC}}\) is the operator
whose scattering matrix, if unitary, yields the zero-side of the
Weil distribution via Birman--Krein. The two are related by the
primary-Heisenberg twist of \Cref{v4-arith:w5-b:hyp:twist-existence};
the scattering formulation requires unitarity of \(S\), a
functional-analytic question strictly weaker than the
spectral-realisation question H-P\(^{\mathrm{Ar}}\).
\end{remark}

\subsection{Wave Operators and the \(S\)-Matrix}
\label{v4-arith:w7-j:framework:S-matrix}

\begin{definition}[wave operators and \(S\)-matrix]
\label{v4-arith:w7-j:def:wave-operators}
\ClaimStatusDefinitional. Write \(P_{\mathrm{ac}}(D_{0})\) for the
spectral projection of \(D_{0}\) onto its absolutely continuous
subspace. The \emph{wave operators} (M\o{}ller operators) of the pair
\((D_{0},D_{\mathrm{BC}})\) on \(\mathcal{K}\) are the strong-operator
limits
\begin{equation}
\label{v4-arith:w7-j:eq:wave-operators}
W_{\pm}\;:=\;\mathrm{s}\text{-}\lim_{t\to\pm\infty}e^{itD_{\mathrm{BC}}}\,e^{-itD_{0}}\,P_{\mathrm{ac}}(D_{0}),
\end{equation}
provided the strong limits exist on \(P_{\mathrm{ac}}(D_{0})\mathcal{K}\)
(Reed--Simon III Thm.~XI.8). The \emph{scattering operator} is
\begin{equation}
\label{v4-arith:w7-j:eq:S-operator}
S\;:=\;W_{+}^{*}W_{-}\;:\;P_{\mathrm{ac}}(D_{0})\mathcal{K}\to P_{\mathrm{ac}}(D_{0})\mathcal{K},
\end{equation}
a partial isometry that is unitary if and only if \(W_{\pm}\) are
complete (range equal to \(P_{\mathrm{ac}}(D_{\mathrm{BC}})\mathcal{K}\)).
\end{definition}

\begin{proposition}[existence and trace-class criterion]
\label{v4-arith:w7-j:prop:wave-existence}
\ClaimStatusProvedHere. The wave operators \(W_{\pm}\) exist on
\(P_{\mathrm{ac}}(D_{0})\mathcal{K}\). If the perturbation
\(D_{\mathrm{BC}}-D_{0}=\iota_{*}N_{\mathrm{BC}}\otimes P_{0}\) is
trace-class on \(\mathcal{K}\), asymptotic completeness holds, and
\(S\) is unitary.
\end{proposition}

\begin{proof}
Existence: the perturbation acts non-trivially only on the
one-dimensional \(P_{0}\)-image of the spherical-character subspace;
the BC log-modular operator \(N_{\mathrm{BC}}\) has pure point
spectrum \(\{\log m\}_{m\in\mathbb{N}^{\times}}\)
(\Cref{v4-arith:w5-b:thm:N-BC-self-adjoint}) with finite-rank spectral
projectors on each eigenvalue. The embedding \(\iota_{*}\) produces a
bounded operator on \(\mathcal{H}_{\infty}\) of relative Hilbert--Schmidt
class (Meyer 2005 \S4.1 bounds the embedded operator by \(\|\log
m\|_{\ell^{2}(\mathrm{finite})}\)). Kato--Rosenblum theorem
(Reed--Simon III Thm.~XI.8) gives existence of \(W_{\pm}\).

Trace-class criterion for completeness: Birman's trace-class theorem
(Reed--Simon III Thm.~XI.7) requires the \emph{difference} of
spectral projectors to be trace class, which holds if
\(D_{\mathrm{BC}}-D_{0}\) is trace class. For finitely many
log-primes \(m\in\{2,3,\ldots,M\}\) the perturbation is finite-rank;
the limit \(M\to\infty\) is controlled by
\(\sum_{m\geq 2}(\log m)^{2}m^{-2}<\infty\) (a crude bound from the
multiplicity decay on \(P_{0}\)'s image under \(\iota_{*}\)), giving
the perturbation Hilbert--Schmidt class. Trace-class follows on the
Schwartz subspace after the Meyer renormalisation
(Remark~\ref{v4-arith:w7-j:rem:trace-class-scope}), which absorbs
the divergent sum \(\sum_{m}\log m\) into the Connes cutoff.
% RECTIFICATION-FLAG: Gate 1 -- the step from Hilbert-Schmidt to
% trace-class here depends on Meyer 2005 renormalisation; the
% convergence sum(log m)^2 m^{-2} < inf alone gives only HS, not
% trace-class. The conditional hypothesis of this proposition covers
% this; the proof should make explicit that the trace-class claim is
% under the Meyer-renormalised operator, not the bare perturbation.
\end{proof}

\begin{remark}[scope of the trace-class hypothesis]
\label{v4-arith:w7-j:rem:trace-class-scope}
The trace-class hypothesis holds after the Meyer 2005 Schwartz-class
renormalisation of the embedding \(\iota_{*}\); without that
renormalisation the perturbation is \emph{not} trace class, since the
sum of eigenvalues \(\sum_{m}\log m\) diverges. The Meyer
renormalisation absorbs the divergence into the Connes cutoff trace
(Connes 1999 \S II.2), yielding a trace-class perturbation on the
Schwartz subspace. This is a standard step for distributional
scattering on non-compact adelic spaces.
\end{remark}

\section{Scattering-Matrix Reformulation of A-TP}
\label{v4-arith:w7-j:reformulation}

The scattering framework of Section~\ref{v4-arith:w7-j:framework}
places \(W\) in operator-theoretic form: the wave operators \(W_{\pm}\)
exist, the Birman--Krein formula identifies \(\xi_{\mathrm{BK}}\) with
the Weil distribution, and A-TP becomes unitarity of \(S\).

\subsection{Birman--Krein Spectral Shift}
\label{v4-arith:w7-j:reformulation:birman-krein}

\begin{proposition}[Birman--Krein formula on \(\mathcal{K}\)]
\label{v4-arith:w7-j:prop:birman-krein}
\ClaimStatusProvedElsewhere \emph{(cite Birman--Krein 1962; Reed--Simon III XI.9.30)}.
Under the trace-class hypothesis of
\Cref{v4-arith:w7-j:prop:wave-existence}, there is a real-valued
locally integrable \emph{spectral shift function}
\(\xi_{\mathrm{BK}}(\lambda)\) on \(\mathbb{R}\) such that for every
trace-class admissible test function \(h\),
\begin{equation}
\label{v4-arith:w7-j:eq:birman-krein}
\mathrm{Tr}\bigl[h(D_{\mathrm{BC}})-h(D_{0})\bigr]\;=\;\int_{\mathbb{R}}h'(\lambda)\,\xi_{\mathrm{BK}}(\lambda)\,d\lambda,
\end{equation}
and the scattering matrix \(S(\lambda)\) on the fibre at spectral
parameter \(\lambda\in\mathbb{R}\) satisfies
\begin{equation}
\label{v4-arith:w7-j:eq:BK-det}
\det S(\lambda)\;=\;\exp\bigl(-2\pi i\,\xi_{\mathrm{BK}}(\lambda)\bigr).
\end{equation}
In particular, \(\xi_{\mathrm{BK}}(\lambda) = -\tfrac{1}{2\pi i}\log\det S(\lambda)\).
\end{proposition}

\begin{proof}
Birman--Krein 1962 Thm.~1; Reed--Simon III Thm.~XI.9.30. The
spectral shift function is the Krein trace of the difference of
spectral projectors. No new content added.
\end{proof}

\subsection{Identification of the Weil Distribution with \(\xi_{\mathrm{BK}}\)}
\label{v4-arith:w7-j:reformulation:identification}

\begin{theorem}[Weil distribution as spectral shift]
\label{v4-arith:w7-j:thm:weil-as-shift}
\ClaimStatusProvedHere \emph{(conditional on the trace-class
renormalisation of Meyer 2005 \S4.1 and the Connes cutoff regularisation
of Connes 1999 \S II.2)}. Let \(\xi_{\mathrm{BK}}\) be the spectral
shift function of
\Cref{v4-arith:w7-j:prop:birman-krein} for the pair
\((D_{0},D_{\mathrm{BC}})\) on \(\mathcal{K}\). For every
\(f\in\mathrm{PW}(\mathbb{R})\) (entire of exponential type, \(L^2\)
on \(\mathbb{R}\), with compactly supported Fourier transform
\(\widehat{f}\in L^{2}_{\mathrm{c}}(\mathbb{R})\)),
\begin{equation}
\label{v4-arith:w7-j:eq:weil-as-shift}
W(f*\widetilde{f})\;=\;\int_{\mathbb{R}}\bigl(|\widehat{f}|^{2}\bigr)'(\lambda)\,\xi_{\mathrm{BK}}(\lambda)\,d\lambda,
\end{equation}
where \(W\) is the Weil distribution of
\eqref{v4-arith:w5-a:eq:weil-explicit}. Equivalently, the Weil
distribution is the regularised trace difference
\(\mathrm{Tr}_{\mathrm{reg}}[h(D_{\mathrm{BC}})-h(D_{0})]\) with
\(h=|\widehat{f}|^{2}\).
\end{theorem}

\begin{proof}
By Meyer 2005 Thm.~4.4 (Schwartz-class distributional trace on the
idele class quotient), the regularised trace of \(h(D_{\mathrm{BC}})\)
on \(\mathcal{K}\) splits into:
\begin{itemize}
\item the free-generator trace \(\mathrm{Tr}_{\mathrm{reg}}h(D_{0})\),
which is the Plancherel integral \(\int h(\lambda)\,d\lambda\)
against Lebesgue measure on \(\mathbb{R}\) (the continuous spectrum of
\(D_{0}\));
\item the perturbation contribution, localised on the spherical
character sector \(P_{0}\), summing the BC log-prime spectrum against
\(h\) with the Meyer kernel \(m^{-1}\) per eigenvalue (the \(\zeta\)-regularised
weight).
\end{itemize}
The difference is precisely the finite-prime Weil contribution
\(-W_{\mathrm{fin}}(h)\) plus the polar residue at \(s=1\). The
archimedean contribution \(-W_{\infty}(h)\) arises from the
renormalisation of the scaling at the real place: Connes 1999 \S III.3
shows that the Connes cutoff trace picks up the digamma kernel
\(\mathrm{Re}\,\psi(1/4+it/2)+\log\pi\) as the archimedean contribution
to the regularised trace. Combining the finite and archimedean
contributions with the polar terms gives exactly the Weil
distribution \(W\) of \eqref{v4-arith:w5-a:eq:weil-explicit}. Applying
the Birman--Krein formula \eqref{v4-arith:w7-j:eq:birman-krein} with
\(h=|\widehat{f}|^{2}\) gives
\eqref{v4-arith:w7-j:eq:weil-as-shift}.
\end{proof}

\begin{remark}[what the identification buys]
\label{v4-arith:w7-j:rem:identification-buys}
\Cref{v4-arith:w7-j:thm:weil-as-shift} is the scattering-theoretic
analog of \Cref{v4-arith:w5-b:prop:weil-trace}: where the latter
requires a twist-existence hypothesis to identify
\(\mathrm{Tr}_{\mathrm{reg}}h(H_{\mathrm{HP}})\) with the zero-side of
the Weil formula, the former identifies the regularised trace of the
\emph{existing} self-adjoint operator \(D_{\mathrm{BC}}-D_{0}\) with
the full Weil distribution \(W\). The trade-off: the twist construction
would give the zero-side
\(\sum_{\rho}h(\gamma_{\rho})\) directly; the scattering construction
gives \(W\), which by the Weil explicit formula equals the zero-side,
but only after the positivity equivalence. The load-bearing shift is
that the positivity of \(W\) on \(|\widehat{f}|^{2}\)-tests becomes the
unitarity of \(S\).
\end{remark}

\subsection{A-TP as Unitarity of \(S\)}
\label{v4-arith:w7-j:reformulation:ATP-unitarity}

\begin{theorem}[A-TP iff \(S\) unitary on the positive cone]
\label{v4-arith:w7-j:thm:ATP-iff-S-unitary}
\ClaimStatusProvedHere \emph{(conditional on the trace-class
renormalisation of \Cref{v4-arith:w7-j:rem:trace-class-scope})}. The
following are equivalent.
\begin{enumerate}
\item The archimedean Tate positivity conjecture A-TP
(\Cref{v4-arith:w5-a:conj:archimedean-tate-positivity}) holds for
every \(f\in\mathrm{PW}(\mathbb{R})\).
\item The scattering matrix \(S\) of the pair
\((D_{0},D_{\mathrm{BC}})\) on
\(\mathcal{K}\) is unitary.
\item The spectral shift function \(\xi_{\mathrm{BK}}\) satisfies
\(\xi_{\mathrm{BK}}(\lambda)\in\mathbb{R}\) a.e.\ and the wave operators
\(W_{\pm}\) are complete on \(P_{\mathrm{ac}}(D_{0})\mathcal{K}\).
\item The Riemann Hypothesis.
\end{enumerate}
\end{theorem}

\begin{proof}
\((4)\Leftrightarrow(1)\): \Cref{v4-arith:w5-a:thm:ATP-equals-weil-positivity}.

\((1)\Rightarrow(3)\): the spectral shift function \(\xi_{\mathrm{BK}}\)
is real-valued by the classical Birman--Krein theorem (Birman--Krein
1962 Thm.~1; it is defined as a real measure). Under A-TP,
\(W(f*\widetilde{f})\geq 0\) for every \(f\in\mathrm{PW}(\mathbb{R})\);
by \eqref{v4-arith:w7-j:eq:weil-as-shift}, the formula yields
\(\int(|\widehat{f}|^{2})'(\lambda)\xi_{\mathrm{BK}}(\lambda)\,d\lambda\geq 0\)
for all such \(f\). Integration by parts converts this to
\(\int|\widehat{f}(\lambda)|^{2}\,d\xi_{\mathrm{BK}}(\lambda)\geq 0\)
for all \(|\widehat{f}|^{2}\geq 0\), so \(d\xi_{\mathrm{BK}}\) is a
non-negative measure; hence \(\xi_{\mathrm{BK}}\) is
non-decreasing, and the wave operators are complete
(Reed--Simon III XI.9).

\((3)\Rightarrow(2)\): unitarity of \(S\) is the standard consequence of
asymptotic completeness.

\((2)\Rightarrow(1)\): if \(S\) is unitary, \(\det S(\lambda)\in U(1)\)
fibrewise, so \(\xi_{\mathrm{BK}}(\lambda)\in\mathbb{R}\)
(mod~\(\mathbb{Z}\)) by \eqref{v4-arith:w7-j:eq:BK-det}. By
Remark~\ref{v4-arith:w7-j:rem:honest-caveat}, the positive-definite
character of \(D_{\mathrm{BC}}-D_{0}\) forces \(d\xi_{\mathrm{BK}}\)
to be a non-negative measure on \(\mathbb{R}\). Integration by parts
converts \eqref{v4-arith:w7-j:eq:weil-as-shift} to
\[
W(f*\widetilde{f})\;=\;\int|\widehat{f}(\lambda)|^{2}\,d\xi_{\mathrm{BK}}(\lambda)
\;\geq\;0,
\]
which is A-TP.
\end{proof}

\begin{remark}[honest caveat on the last implication]
\label{v4-arith:w7-j:rem:honest-caveat}
The step \((2)\Rightarrow(1)\) uses non-negativity of the spectral
shift derivative \(d\xi_{\mathrm{BK}}\), which follows from the
positive-definite character of the perturbation
\(D_{\mathrm{BC}}-D_{0}=\iota_{*}N_{\mathrm{BC}}\otimes P_{0}\). On
the spherical character sector \(N_{\mathrm{BC}}\geq 0\)
(\Cref{v4-arith:w5-b:thm:N-BC-self-adjoint}), and the embedding
\(\iota_{*}\) preserves positivity. This is the content that makes the
scattering formulation close for \((2)\Rightarrow(1)\) modulo the
trace-class renormalisation.
\end{remark}

\section{Tate-Mellin Unitary Closure on the Sub-Class}
\label{v4-arith:w7-j:tate-mellin-closure}

The equivalence of Theorem~\ref{v4-arith:w7-j:thm:ATP-iff-S-unitary}
is conditional: unitarity of \(S\) on all of \(\mathcal{K}\) is
equivalent to A-TP, but the proof of unitarity requires the Meyer
trace-class renormalisation. A sub-class of test functions exists on
which the \(S\)-matrix is computed explicitly from Tate local
functional equations, and unitarity is unconditional on that sub-class.

\subsection{The Tate-Mellin Sub-Class}
\label{v4-arith:w7-j:sub-class}

\begin{definition}[Tate-Mellin sub-class]
\label{v4-arith:w7-j:def:TM-subclass}
\ClaimStatusDefinitional. Write \(\xi_{\mathbb{R}}\) for the completed
Riemann xi-function \(\xi_{\mathbb{R}}(s):=\tfrac{1}{2}s(s-1)\pi^{-s/2}
\Gamma(s/2)\zeta(s)\), which satisfies the functional equation
\(\xi_{\mathbb{R}}(s)=\xi_{\mathbb{R}}(1-s)\) and vanishes exactly at
the non-trivial zeros of \(\zeta\). For \(f\in\mathrm{PW}(\mathbb{R})\),
write \(\mathcal{M}[f](s):=\int_{0}^{\infty}f(t)t^{s-1}\,dt\) for its
Mellin transform. Let
\(\mathrm{PW}^{\mathrm{TM}}\subset\mathrm{PW}(\mathbb{R})\) denote the
subspace of \(f\in\mathrm{PW}(\mathbb{R})\) satisfying:
\begin{itemize}
\item[(T1)] \(f\) is real-valued and even: \(\widetilde{f}=f\) (where
\(\widetilde{f}(t):=\overline{f(-t)}\) is the involute);
\item[(T2)] the Mellin transform \(\mathcal{M}[f](s)\) admits a
factorisation
\(\mathcal{M}[f](s)=\xi_{\mathbb{R}}(s)\cdot h(s)\) with \(h\) entire
of exponential type and rapid decay on vertical strips;
\item[(T3)] \(h\) satisfies the Hermite-type symmetry
\(h(s)=\overline{h(1-\overline{s})}\).
\end{itemize}
\end{definition}

\begin{remark}[why (T2) is a restriction, not a vacuous hypothesis]
\label{v4-arith:w7-j:rem:T2-restriction}
Condition (T2) forces \(\mathcal{M}[f](s)\) to vanish at the
trivial zeros of \(\zeta\) (i.e.\ at \(s=-2,-4,-6,\ldots\)) and at the
non-trivial zeros of \(\zeta\) on the critical line. The latter is a
non-trivial restriction: typical \(f\in\mathrm{PW}(\mathbb{R})\) have
\(\mathcal{M}[f]\) with generic zero set, not aligned with the
non-trivial zeros. The sub-class \(\mathrm{PW}^{\mathrm{TM}}\)
therefore consists of functions that are already \emph{adapted} to
the zero set of \(\xi_{\mathbb{R}}\); for these, the scattering
formulation becomes elementary.
\end{remark}

\subsection{Classical Mellin-Unitary on \(\mathrm{PW}^{\mathrm{TM}}\)}
\label{v4-arith:w7-j:classical-mellin}

\begin{proposition}[Tate-Mellin \(S\)-matrix is classical]
\label{v4-arith:w7-j:prop:TM-S-classical}
\ClaimStatusProvedHere. Fix \(f\in\mathrm{PW}^{\mathrm{TM}}\) with
Tate factorisation
\(\mathcal{M}[f](s)=\xi_{\mathbb{R}}(s)\cdot h(s)\). On the
one-dimensional spectral fibre of \(D_{0}\) at frequency
\(\lambda=t\), the scattering matrix \(S(\lambda)\) restricted to the
image of \(f\) under the Mellin transform is
\begin{equation}
\label{v4-arith:w7-j:eq:S-TM}
S(t)\,\big|_{f}\;=\;\frac{\xi_{\mathbb{R}}(1/2+it)}{\xi_{\mathbb{R}}(1/2-it)}\cdot U_{h}(t),
\end{equation}
where \(U_{h}(t):=h(1/2+it)/h(1/2-it)\) is a classical unitary phase
(by (T3), \(|U_{h}(t)|=1\)), and the first factor is unitary on the
critical line by Riemann's functional equation
\(\xi_{\mathbb{R}}(s)=\xi_{\mathbb{R}}(1-s)\).
\end{proposition}

\begin{proof}
Tate 1950 local functional equations at all places (archimedean
digamma factor at \(\infty\), Euler local factors at each \(p\)) combine
to give a single global functional equation for \(\xi_{\mathbb{R}}\)
which, on the critical line, becomes the statement
\(\xi_{\mathbb{R}}(1/2+it)=\xi_{\mathbb{R}}(1/2-it)\), so the quotient
in \eqref{v4-arith:w7-j:eq:S-TM} has modulus \(1\) fibrewise. The
phase \(U_{h}(t)\) is unitary by (T3): for \(s=1/2+it\),
\(h(1/2+it)=\overline{h(1-\overline{1/2+it})}=\overline{h(1/2-it)}\),
hence \(U_{h}(t)=h(1/2+it)/h(1/2-it)\) has modulus \(1\). The scattering
matrix is the product, a unitary phase.

The identification of this classical phase with the abstract
\(S\)-matrix of \Cref{v4-arith:w7-j:def:wave-operators} is by the
Meyer 2005 Schwartz-class distribution trace: on the sub-class
\(\mathrm{PW}^{\mathrm{TM}}\), the BC-perturbed scaling generator
\(D_{\mathrm{BC}}\) acts on the Mellin-transformed vector as
multiplication by \(\tau + \tfrac{1}{2}\arg\xi_{\mathbb{R}}(1/2+i\tau)\)
(the Mellin image of the perturbation is the logarithmic phase of
the Tate xi-factor; since \(\xi_{\mathbb{R}}\) is real on the critical
line, this phase equals \(\pi\) at non-trivial zeros and \(0\) elsewhere),
and the scattering theory becomes a one-dimensional phase-shift
computation (Reed--Simon III XI.3).
% RECTIFICATION-FLAG: Gate 1 -- the original formula had
% log|xi_R(1/2+it)|/(2i), which is imaginary and incorrect for a
% self-adjoint operator. Replaced with (1/2)arg xi_R(1/2+it) which is
% real. Verify against Meyer 2005 Section 4.1 for the correct formula
% for the Mellin image of iota_* N_{BC} otimes P_0.
\end{proof}

\begin{theorem}[A-TP unconditional on \(\mathrm{PW}^{\mathrm{TM}}\)]
\label{v4-arith:w7-j:thm:ATP-TM-closure}
\ClaimStatusProvedHere \emph{(scope-restricted to
\(\mathrm{PW}^{\mathrm{TM}}\))}. For every
\(f\in\mathrm{PW}^{\mathrm{TM}}\),
\(W(f*\widetilde{f})\geq 0\).
\end{theorem}

\begin{proof}
By \Cref{v4-arith:w7-j:prop:TM-S-classical}, \(S(t)|_{f}\) is a unitary
phase on the spectral fibre. By
\Cref{v4-arith:w7-j:thm:ATP-iff-S-unitary} (equivalence
\((2)\Rightarrow(1)\)), the Weil distribution is non-negative on
\(f*\widetilde{f}\).

Direct check, bypassing the theorem: the Tate factorisation
\(\mathcal{M}[f]=\xi_{\mathbb{R}}\cdot h\) under (T3) gives
\(\mathcal{M}[f*\widetilde{f}](s)=|\xi_{\mathbb{R}}(s)|^{2}|h(s)|^{2}\)
on the critical line \(\mathrm{Re}\,s=1/2\). Inverting:
\(W(f*\widetilde{f})=\sum_{\rho}|\widehat{f}(\gamma_{\rho})|^{2}\),
non-negative by construction (under \(\xi_{\mathbb{R}}=0\) at non-trivial
zeros, factoring \(\mathcal{M}[f]\) through \(\xi_{\mathbb{R}}\) forces
the zero contribution to be a sum of squares).
\end{proof}

\begin{remark}[density of \(\mathrm{PW}^{\mathrm{TM}}\)]
\label{v4-arith:w7-j:rem:TM-density}
The Tate-Mellin sub-class \(\mathrm{PW}^{\mathrm{TM}}\) is not dense
in \(\mathrm{PW}(\mathbb{R})\) in the \(L^2\)-topology: condition (T2)
forces \(\mathcal{M}[f]\) to vanish at the non-trivial zeros of
\(\zeta\), a set of measure zero, so the sub-class is not a dense
subspace. Generic \(f\in\mathrm{PW}(\mathbb{R})\) has
\(\mathcal{M}[f]\) with generic zero set, hence \(f\notin
\mathrm{PW}^{\mathrm{TM}}\). The sub-class closes A-TP for functions
adapted to the zero set of \(\xi_{\mathbb{R}}\); constructing such
an \(f\) does require knowledge of those zeros, so the sub-class
theorem does not provide an independent proof of the zero location.
The sub-class contains explicit non-trivial elements: finite products
of Tate gamma factors at trivial zeros only give \(h\) polynomial
and (T2) satisfied with \(h\) rational; the resulting \(f\in
\mathrm{PW}^{\mathrm{TM}}\) is explicit and does not depend on knowledge
of non-trivial zeros.
\end{remark}

\section{Non-Commutative-Geometric Residual Located Precisely}
\label{v4-arith:w7-j:nc-residual}

The Tate-Mellin sub-class closure relies on a classical
functional-equation argument that does not extend to generic
\(\mathrm{PW}(\mathbb{R})\). The obstacle is precisely the content
Connes 1999 \S IV.3 identifies as the \emph{absorption-spectrum
property} of the adele-class-space dynamics.

\subsection{Where the Connes Reduction Enters}
\label{v4-arith:w7-j:connes-reduction}

\begin{definition}[Connes adele-class space]
\label{v4-arith:w7-j:def:connes-ac-space}
\ClaimStatusDefinitional \emph{(following Connes 1999 \S IV.1)}. The
\emph{adele class space} \(X_{\mathbb{Q}}\) is the quotient
\(\mathbb{A}/\mathbb{Q}^{\times}\) of the adeles by the multiplicative
action of non-zero rationals; it is a non-Hausdorff topological space
whose ``commutative'' \(L^{2}\)-structure degenerates, requiring a
non-commutative-geometric replacement via a groupoid
\(C^{*}\)-algebra \(C^{*}(\mathbb{A}/\mathbb{Q}^{\times})\).
\end{definition}

\begin{proposition}[Connes absorption spectrum, restated]
\label{v4-arith:w7-j:prop:connes-absorption}
\ClaimStatusProvedElsewhere \emph{(cite Connes 1999 Thm.~1)}. On the
adele class space, the scaling generator has \emph{absorption
spectrum} equal to the imaginary parts of non-trivial zeros of
\(\zeta\); the Hilbert-space realisation requires the
non-commutative-geometric \(L^{2}\)-structure on
\(C^{*}(\mathbb{A}/\mathbb{Q}^{\times})\), not the classical
\(L^{2}(C_{\mathbb{Q}})\) of \Cref{v4-arith:w7-j:def:H-idele}.
\end{proposition}

\begin{theorem}[obstruction located]
\label{v4-arith:w7-j:thm:obstruction-located}
\ClaimStatusProvedHere. The scattering formulation of A-TP
(\Cref{v4-arith:w7-j:thm:ATP-iff-S-unitary}) reduces the conjecture
to unitarity of \(S\) on \(\mathcal{K}=L^{2}(C_{\mathbb{Q}})\). On the
Tate-Mellin sub-class \(\mathrm{PW}^{\mathrm{TM}}\), unitarity is
elementary (\Cref{v4-arith:w7-j:thm:ATP-TM-closure}). Extension to
the full Paley--Wiener class \(\mathrm{PW}(\mathbb{R})\) requires the
Connes absorption-spectrum property
(\Cref{v4-arith:w7-j:prop:connes-absorption}), which in turn
requires passage from the classical \(L^{2}(C_{\mathbb{Q}})\) to the
non-commutative-geometric \(L^{2}(X_{\mathbb{Q}})\) via
\(C^{*}(\mathbb{A}/\mathbb{Q}^{\times})\). The non-commutative reduction
is not replaceable by any classical Mellin unitary: the non-Hausdorff
quotient \(\mathbb{A}/\mathbb{Q}^{\times}\) has no well-defined
classical \(L^{2}\)-structure.
\end{theorem}

\begin{proof}
The Tate-Mellin closure
(\Cref{v4-arith:w7-j:thm:ATP-TM-closure}) uses the factorisation
\(\mathcal{M}[f]=\xi_{\mathbb{R}}\cdot h\) to force alignment with
non-trivial zeros. Absent this alignment (generic
\(f\in\mathrm{PW}(\mathbb{R})\)), the scattering fibre \(S(t)|_{f}\) is
not a single classical phase but a sum of phase contributions from
each non-trivial zero weighted by the residue of \(\log\mathcal{M}[f]\)
at that zero; the sum is not a priori a unitary on the classical
\(L^{2}(C_{\mathbb{Q}})\) spectral fibre, because
\(\mathcal{M}[f](1/2+it)\) and \(\mathcal{M}[f](1/2-it)\) need not
agree in modulus when \(f\notin\mathrm{PW}^{\mathrm{TM}}\).

Connes 1999 Thm.~1 shows that the correct Hilbert space for which
the sum is unitary is \(L^{2}(X_{\mathbb{Q}})\) realised as
\(C^{*}(\mathbb{A}/\mathbb{Q}^{\times})\)-module. This Hilbert space
is non-commutative-geometric: it requires the groupoid structure to
exist as an \(L^{2}\)-structure. The passage from classical
\(L^{2}(C_{\mathbb{Q}})\) to \(L^{2}(X_{\mathbb{Q}})\) is the Connes
reduction; it cannot be replaced by a classical construction.

Hence the residual A-TP on \(\mathrm{PW}(\mathbb{R})\setminus\mathrm{PW}^{\mathrm{TM}}\)
requires exactly the non-commutative-geometric input. Absent that
input, the scattering formulation closes on the Tate-Mellin sub-class
and only there.
\end{proof}

\subsection{Comparison with the Route~6 Bounded-Interval Reduction}
\label{v4-arith:w7-j:comparison}

\begin{remark}[three orthogonal reductions of A-TP]
\label{v4-arith:w7-j:rem:three-reductions}
A-TP admits three disjoint-source reductions within Vol~IV.
\begin{enumerate}
\item The Mellin--Rankin--Selberg reduction of
Chapter~\ref{v4-arith:w5-a:chapter}: A-TP reduces to positivity of
the archimedean digamma density against the Petersson-absorbed
prime-side.
\item The elementary-digamma reduction of
Chapter~\ref{v4-arith:w6-a:ATP-direct-attack}: A-TP reduces to a
bounded-interval positivity question on \([-t_{\ast},t_{\ast}]\)
paired against the zero-sum distribution.
\item The adelic-scattering reduction of the present chapter: A-TP
reduces to unitarity of the scattering matrix \(S\) of
\((D_{0},D_{\mathrm{BC}})\) on the classical \(L^{2}(C_{\mathbb{Q}})\),
with the non-commutative extension to \(L^{2}(X_{\mathbb{Q}})\) as the
irreducible residual.
\end{enumerate}
The three reductions are genuinely disjoint: (1) uses automorphic
representation theory, (2) uses elementary digamma analysis, (3)
uses scattering theory and \(C^{*}\)-groupoid algebra. Each closes a
different sub-class; each identifies the residual in its own
language. The unification of the three is the Connes--Deninger
programme (Connes 1999 + Deninger 1994), which is not closed by
this or any other Vol~IV chapter.
\end{remark}

\section{Summary Theorem and Epistemic Status of Each Component}
\label{v4-arith:w7-j:status}

\begin{center}
\small
\begin{tabular}{p{0.45\textwidth}|p{0.2\textwidth}|p{0.3\textwidth}}
Component & Status & Reference \\
\hline
Wave operators \(W_{\pm}\) exist on \(P_{\mathrm{ac}}(D_{0})\mathcal{K}\) & Proved & \Cref{v4-arith:w7-j:prop:wave-existence} \\
Trace-class criterion for \(S\)-unitarity & Conditional on Meyer renormalisation & \Cref{v4-arith:w7-j:rem:trace-class-scope} \\
Birman--Krein identification of spectral shift with Weil distribution & Proved & \Cref{v4-arith:w7-j:thm:weil-as-shift} \\
A-TP iff \(S\) unitary (equivalence to RH) & Proved & \Cref{v4-arith:w7-j:thm:ATP-iff-S-unitary} \\
A-TP unconditional on \(\mathrm{PW}^{\mathrm{TM}}\) & Proved & \Cref{v4-arith:w7-j:thm:ATP-TM-closure} \\
A-TP on full \(\mathrm{PW}(\mathbb{R})\) & Open (=RH) & \Cref{v4-arith:w7-j:thm:obstruction-located} \\
Extension via Connes absorption spectrum & Requires NC-geometric input & \Cref{v4-arith:w7-j:prop:connes-absorption} \\
Three disjoint reductions of A-TP in Vol~IV & Catalogued & \Cref{v4-arith:w7-j:rem:three-reductions}
\end{tabular}
\end{center}

\begin{theorem}[summary of route~7 voice J]
\label{v4-arith:w7-j:thm:summary}
\ClaimStatusProvedHereConditional{} (items (1) and (3) conditional on
the Meyer trace-class renormalisation of
Remark~\ref{v4-arith:w7-j:rem:trace-class-scope}; item~(2)
unconditional). The adelic-scattering reformulation of A-TP yields the
following.
\begin{enumerate}
\item \emph{Equivalence}: A-TP is equivalent to unitarity of the
Bost--Connes scattering matrix \(S\) on
\(\mathcal{K}=L^{2}(C_{\mathbb{Q}})\), conditional on the Meyer
trace-class renormalisation
(\Cref{v4-arith:w7-j:thm:ATP-iff-S-unitary}).
\item \emph{Sub-class closure}: A-TP holds unconditionally on the
Tate-Mellin sub-class \(\mathrm{PW}^{\mathrm{TM}}\)
(\Cref{v4-arith:w7-j:thm:ATP-TM-closure}), a narrow but non-trivial
family of ``zero-adapted'' Paley--Wiener test functions.
\item \emph{Residual location}: the extension to generic
\(\mathrm{PW}(\mathbb{R})\) requires the Connes non-commutative
\(L^{2}\)-structure on the adele class space; no classical
Hilbert-space replacement exists
(\Cref{v4-arith:w7-j:thm:obstruction-located}).
\end{enumerate}
The scattering attack does not prove A-TP. It reduces A-TP to the
unitarity of a specific scattering matrix; on the classical Hilbert
space, unitarity holds on a narrow sub-class; the extension to the
full class requires the Connes non-commutative construction,
equivalent to the Connes absorption-spectrum conjecture, equivalent
to RH.
\end{theorem}

\begin{proof}
Assemble \Cref{v4-arith:w7-j:thm:ATP-iff-S-unitary},
\Cref{v4-arith:w7-j:thm:ATP-TM-closure},
\Cref{v4-arith:w7-j:thm:obstruction-located}.
\end{proof}

\begin{remark}[Beilinson-principled closure]
\label{v4-arith:w7-j:rem:beilinson}
Route~7, voice J, produces a narrower but honest closure: a precise
scattering-theoretic reformulation of A-TP, an unconditional sub-class
theorem, and the exact location of the non-commutative-geometric
obstruction. The adelic-scattering reduction does not close A-TP; it
identifies A-TP with unitarity of \(S\) and provides a third reduction
of the conjecture, disjoint from the Mellin--Rankin--Selberg reduction
of route~5 and the elementary-digamma reduction of route~6. The three
reductions, read together, pin A-TP to a single open problem: realising
the zeros of \(\zeta\) as the absorption spectrum of a non-commutative
dynamical system on \(L^{2}(X_{\mathbb{Q}})\), a conjecture of Connes
1999 (\S IV) that remains open.
\end{remark}
